% ===========================================================================
%  E09 — Inteligência artificial com limites
% ===========================================================================
\chapter{E09 \textperiodcentered{} Inteligência artificial com limites}
\label{ep:e09}

\begin{fichaepico}
\begin{campos}
  \item[Objetivo] Usar IA onde ela poupa trabalho do professor ou enriquece a
    missão --- sem que ela decida nota, conheça a criança ou seja condição
    para a aula acontecer. A única vez em que ela dá um nível é a exceção
    declarada da A Seleção, e quem decide continua sendo o professor
    (\id{H09.3.1}).
  \item[Para quem] Professor e Estudante (quem usa); rede (quem precisa
    confiar).
  \item[Valor] IA na escola pública só é aceitável com garantia e com custo
    previsível. Este épico é essa garantia.
  \item[Estado] \parcial{} --- as garantias estão prontas; duas missões ainda
    usam IA fora do limite diário.
\end{campos}
\end{fichaepico}

\begin{tabularx}{\linewidth}{@{}L{22mm}X X@{}}
\toprule
\headrow \thd{Onde} & \thd{O que a IA faz} & \thd{O que ela nunca faz}\\
\midrule
Preparo da aula & Adapta legendas, mensagens, casos e mídia ao tema do
  professor (Foto Conta, Manifesto, Instanooria, Quem mandou isso?, A Seleção
  --- o tema, os 12 posts, as imagens e as três falas do podcast que dependem
  do tema); gera a voz do podcast da A Seleção. &
  Chegar à turma sem aprovação.\\
\zebra Durante a missão & Responde ao estudante no unooGPT; reações de
  personagens no Instanooria e na Vila; na A Seleção, lê a legenda publicada e
  dá o nível dela numa habilidade da BNCC. & Dar nota, mostrar nível à criança
  ou receber o nome dela. O nível da legenda da A Seleção é a exceção
  declarada: só o professor o vê, e ele confere e ajusta.\\
Planejamento & Gera o plano anual e o das semanas. & Decidir quantas aulas
  cabem ou inventar habilidade da BNCC.\\
\zebra Bem-estar & Escreve o texto em volta (abertura, acolhimento). &
  Reescrever o questionário, ou falar quando há sinal de risco.\\
\bottomrule
\end{tabularx}

\begin{tabularx}{\linewidth}{@{}L{16mm}X L{20mm}@{}}
\toprule
\headrow \thd{ID} & \thd{Feature} & \thd{Estado}\\
\midrule
\id{F09.1} & Nada gerado chega à turma sem o professor & \pronto\\
\zebra \id{F09.2} & Nenhuma aula depende da IA & \pronto\\
\id{F09.3} & A IA não dá nota nem conhece a criança & \pronto\\
\zebra \id{F09.4} & Limite diário por perfil & \parcial\\
\bottomrule
\end{tabularx}

% ---------------------------------------------------------------------------
\feature{F09.1}{Nada gerado chega à turma sem o professor}{\pronto}

\begin{historia}{H09.1.1}{Aprovar o que a IA preparou}{\pronto}
\hu{Professor}{revisar e aprovar tudo o que a IA preparou para a minha
aula}{responder pelo que a turma vê}
\aceite
\begin{itemize}
  \item Nas missões com preparo por IA, nada chega aos estudantes sem
    \enquote{Aprovar} (ou \enquote{Aprovar e publicar}); no Instanooria, sem
    \enquote{Ativar}.
  \item Posso revisar item a item, editar, regerar ou excluir.
  \item A tela diz: \enquote{nada gerado por inteligência artificial chega aos
    estudantes sem aprovação}.
\end{itemize}
\end{historia}

% ---------------------------------------------------------------------------
\feature{F09.2}{Nenhuma aula depende da IA}{\pronto}

\begin{historia}{H09.2.1}{Dar a aula mesmo se a IA falhar}{\pronto}
\hu{Professor}{que a aula aconteça mesmo se a IA estiver fora do ar ou sem
saldo}{não perder a aula por um problema que não é meu}
\aceite
\begin{itemize}
  \item Sem IA ou sem saldo, entra o material pronto da Unooria, completo
    (Foto Conta, Manifesto, Quem mandou isso?, A Seleção).
  \item Na A Seleção, sem imagem gerada o post mostra a descrição da cena; sem
    a voz do podcast, a transcrição continua na tela; e, se a leitura da
    legenda falhar, nada é marcado e a enquete continua medindo a habilidade.
  \item No Instanooria, se as reações ao vivo falharem, simplesmente ninguém
    responde --- a missão segue.
  \item No bem-estar, entra um texto fixo e o acompanhamento nunca deixa de
    abrir.
  \item O estudante nunca espera uma geração, exceto na conversa do unooGPT,
    que é a própria missão.
\end{itemize}
\end{historia}

% ---------------------------------------------------------------------------
\feature{F09.3}{A IA não dá nota nem conhece a criança}{\pronto}

\begin{historia}{H09.3.1}{Nota é da plataforma, não da IA}{\pronto}
\hu{rede}{que nenhuma nota, correção ou classificação venha da IA, a não ser
por uma exceção declarada e cercada de garantias}{que o resultado seja justo e
explicável}
\aceite
\begin{itemize}
  \item \textbf{A regra:} gabarito, correção e classificação são calculados
    pela plataforma com regras fixas, e a IA só redige textos.
  \item \textbf{A exceção declarada, e só ela:} na A Seleção, a IA lê a
    legenda que o estudante publica nos Stories e dá o nível dela na EF07CO11
    (atingiu, em processo ou não atingiu), porque só uma leitura do texto diz se
    a legenda explica o critério da escolha (\id{H06.12.5}).
  \item As garantias da exceção: a IA trabalha no modo mais previsível que tem,
    com uma tentativa só, de poucos segundos, e a leitura conta no limite
    diário do estudante. Se qualquer coisa falhar, nada é marcado e a
    habilidade fica \enquote{Não avaliado}; nada do que a IA escreve é guardado.
  \item A criança não vê o nível. O professor o vê no painel de habilidades,
    confere e ajusta, e a primeira enquete da Amara mede a mesma habilidade
    (vale a melhor das duas).
  \item Nas outras missões vale a regra: nenhuma usa a IA para dar nível.
  \item O texto da IA é conferido antes de aparecer (tamanho, termos
    proibidos); se não passar, é refeito ou substituído pelo texto pronto.
\end{itemize}
\end{historia}

\begin{historia}{H09.3.2}{A IA não sabe quem é a criança}{\pronto}
\hu{família do estudante}{que a IA não receba dados que identifiquem a
criança}{que o uso de IA não exponha o meu filho}
\aceite
\begin{itemize}
  \item A IA recebe no máximo a série da turma (\enquote{4º ano}), nunca nome,
    escola ou pontuação.
  \item As instruções da IA nunca vão para o aparelho do estudante.
\end{itemize}
\pendencia{a IA processa fora do Brasil. Isso é aceito, mas precisa constar
no contrato e no termo de tratamento de dados de cada rede (ver o
Capítulo~\ref{cap:backlog}).}
\end{historia}

% ---------------------------------------------------------------------------
\feature{F09.4}{Limite diário por perfil}{\parcial}

\begin{historia}{H09.4.1}{Saber que cheguei ao limite do dia}{\pronto}
\hu{Estudante ou Professor}{uma mensagem clara quando atingi o limite de IA do
dia}{entender o que aconteceu e quando volta}
\aceite
\begin{itemize}
  \item O saldo é conferido \emph{antes} de usar a IA; sem saldo, aparece
    \enquote{limite do dia atingido} e nada é consumido.
  \item O limite é por pessoa e por perfil (estudante ou professor), definido
    pela Unooria (\id{H10.4.1}).
\end{itemize}
\pendencia{o contador zera à meia-noite no horário universal, que no Brasil é
21h. Ver a decisão \id{D04}.}
\end{historia}

\begin{historia}{H09.4.2}{Toda IA conta no limite}{\afazer}
\hu{Administrador Unooria}{que todo uso de IA da plataforma conte no limite
diário}{que o custo de cada rede seja previsível}
\aceiteproposta
\begin{itemize}
  \item A geração e as reações do Instanooria e as reações da Vila do Lume
    passam a contar no limite, como as demais.
  \item O limite do professor comporta a preparação de uma atividade completa
    do Instanooria.
\end{itemize}
\end{historia}

\begin{historia}{H09.4.3}{Ser avisado antes do limite estourar}{\afazer}
\hu{Administrador Unooria}{ser avisado quando o consumo de IA de uma rede se
aproximar do previsto}{ajustar antes que alguém fique sem IA no meio da aula}
\aceiteproposta
\begin{itemize}
  \item Aviso quando o consumo do dia chega a 80\% do previsto para a rede.
\end{itemize}
\end{historia}
